\section{Introduction to Drum
Synthesis}\label{introduction-to-drum-synthesis}

This chapter introduces the basics of drum synthesis, focused on the
classic analogue sounds as a starting point for further exploration.

\subsection{Kicks}\label{kicks}

The basis of nearly every kick sound is an oscillator combined with a
decaying envelope modulating its pitch. Voices 1--3 are ideal for this.

\subsubsection{Oscillator}\label{oscillator}

A good starting point is a sine oscillator with a pitch around 30--35.

\subsubsection{Amplitude Envelope}\label{amplitude-envelope}

A kick starts at maximum amplitude, so set attack to zero. The decay
controls the length of the kick --- a decay in the range 20--30 with a
slope around 25 (exponential) gives a nice average kick. For booming
808-style kicks, increase the decay further.

\subsubsection{Pitch Envelope}\label{pitch-envelope}

Set a decay and slope similar to the amplitude envelope as a starting
point. All three parameters --- decay, slope, and modulation amount ---
have a large impact on the character of the kick.

\subsubsection{Click}\label{click}

The attack transient is an important part of a kick sound. Try different
transient generator settings for various clicks and pops, or slightly
increase the amplitude attack for softer bass drums.

\subsubsection{Tips and Tricks}\label{tips-and-tricks}

\begin{itemize}
\tightlist
\item
  You can further shape the attack using the FM oscillator.
\item
  For additional click in the attack, enhance it with a peak filter.
\end{itemize}

\subsection{Snare}\label{snare}

The snare is a 2-component sound: the tonal body and the noisy rattle.
Voice 4 is designed with both parts in mind --- an oscillator for the
tonal part and a noise generator with filter for the noise part. Only
the noise is routed through the filter, since in most cases you want
highpass-filtered noise with an unfiltered drum body.

For most snare sounds the noise should be louder than the tonal part ---
try setting the oscillator page mix parameter to around 100.

\subsubsection{Tonal Part}\label{tonal-part}

Similar to a kick but less pronounced. A snare does not need a loud
click or deep pitch modulation, and its oscillator frequency should be
higher than a kick since it represents a smaller drum.

\subsubsection{Noise Part}\label{noise-part}

Highpass-filtered noise with a moderate resonance setting is sufficient
for good results.

\subsubsection{Amplitude Envelope}\label{amplitude-envelope-1}

Set attack to zero. For the decay:

\begin{itemize}
\tightlist
\item
  Longer decay times with a very exponential slope (as low as 5--10)
  produce a short hit with a nicely decaying noise tail --- the tail
  acts like a room reverb, giving a more natural result.
\item
  Shorter decay times with a linear or exponential slope give a very
  dry, direct snare.
\end{itemize}

\subsubsection{Tips and Tricks}\label{tips-and-tricks-1}

\begin{itemize}
\tightlist
\item
  Not every snare needs both tone and noise. Try using only the tonal
  part with a high pitch modulation envelope to get the classic
  Kraftwerk electronic zap.
\end{itemize}

\subsection{Clap}\label{clap}

A clap sound is noise with a characteristically fuzzy attack ---
simulating multiple hands clapping with slight timing variations.

\subsubsection{Amplitude Envelope}\label{amplitude-envelope-2}

Use the repeat feature to simulate this behaviour. 3--4 repeats works
well for a classic clap. The timing differences between claps are very
small, so use a short repeat time (attack value below 10). This produces
a short burst at the start. For the decay, use an exponential slope with
a value below 10.

\subsubsection{Oscillator}\label{oscillator-1}

A single white noise oscillator set to its maximum frequency.

\subsubsection{Filter}\label{filter}

A bandpass filter with cutoff 60--80 and high resonance gives good
results.

\subsection{Hihat}\label{hihat}

\subsubsection{Oscillator}\label{oscillator-2}

Hihats need a complex metallic noise spectrum. Use all 3 oscillators
with high frequencies and modulation amounts to build a spectrum you
like. A simple white noise oscillator can also work well.

\subsubsection{Amplitude Envelope}\label{amplitude-envelope-3}

Set 2 decay times --- one for the closed hat, one for the open hat. Both
should be short, but the open hihat decay should be longer than the
closed. Use a very exponential slope (below 10).

\subsubsection{Filter}\label{filter-1}

Use a highpass filter to remove all low frequencies. A high cutoff
around 120 works well.

\subsection{Cymbal Ride}\label{cymbal-ride}

Cymbals are similar to hihats --- the difference is in the envelope and
filter settings.

\subsubsection{Oscillator}\label{oscillator-3}

Oscillator settings can be similar to those of the hihat.

\subsubsection{Amplitude Envelope}\label{amplitude-envelope-4}

The only real difference from the hihat is a much longer decay time. Use
a strong exponential slope. Also consider lowering the voice volume, as
cymbal sounds can be very loud.

\subsubsection{Filter}\label{filter-2}

A bandpass filter with high resonance and a cutoff above 110 gives good
results.

\subsubsection{LFO}\label{lfo}

Cymbals benefit greatly from LFO modulation on the filter frequency. Use
a low modulation frequency and a sine wave with a gentle modulation
depth --- this ensures successive strikes do not sound identical and the
frequency spectrum shifts slightly over time.

\subsection{Bells}\label{bells}

\subsubsection{Realistic Bells}\label{realistic-bells}

Realistic bells are best achieved with FM oscillators. A complex
metallic spectrum combined with an exponential decay amplitude envelope
and a bandpass filter produces different bell types. The sample rate
reducer is very useful for adding a more metallic character.

\subsubsection{808-Style Cowbell}\label{style-cowbell}

The 808 cowbell uses 2 detuned rectangle oscillators through a bandpass
filter.

\textbf{Recommended voice:} Drum voice 1--3

\textbf{FM Page:}

\begin{itemize}
\tightlist
\item
  Set mode to Mix (main and FM oscillators are mixed).
\item
  Set amount to 63 (50/50 mix ratio).
\item
  Select rectangle waveform.
\item
  Set frequency to approximately 78.
\end{itemize}

\textbf{OSC Page:}

\begin{itemize}
\tightlist
\item
  Select rectangle wave.
\item
  Set frequency to 71.
\end{itemize}

\textbf{Amplitude Envelope:}

\begin{itemize}
\tightlist
\item
  Attack: 0
\item
  Slope: very exponential (around 2--5).
\item
  Decay: 25--50 depending on slope setting.
\end{itemize}

\textbf{Filter:}

\begin{itemize}
\tightlist
\item
  Bandpass filter.
\item
  Medium resonance around 70.
\item
  Cutoff frequency around 92.
\end{itemize}

\subsection{More Information About Drum Sound
Design}\label{more-information-about-drum-sound-design}

Two highly recommended resources:

\begin{itemize}
\item
  The \emph{Sound on Sound} Synth Secrets series contains excellent
  articles on different drum sounds:
  \url{http://www.soundonsound.com/search?url=\%2Fsearch&Keyword=\%22synth+secrets\%22}
\item
  The Waldorf Attack synthesizer owners manual contains detailed how-tos
  on drum synthesis:
  \url{http://waldorf.electro-music.com/attack/docs/attackdrumsounds.pdf}
\end{itemize}
